% Source: https://artofproblemsolving.com/community/c3h1878310_question
% Stable identity: aops-topic-1878310-q01
The key fact is $\left(\sqrt2-1\right)\left(\sqrt2+1\right)=2-1=1$. Let $u=\sqrt{\sqrt2-1}>0$, so $u^2=\sqrt2-1$ and $\sqrt2+1=\dfrac{1}{u^2}$, hence $\sqrt{1+\sqrt2}=\dfrac{1}{u}$.

(a) The first product is $u\cdot u^2=u^3$ and the second is $u^2\cdot\dfrac{1}{u}=u$. They are equal only if $u^3=u$, i.e.\ $u=1$, i.e.\ $\sqrt2-1=1$, which is false. So \textbf{no}, they are not equal: numerically $u^3\approx0.267$ and $u\approx0.644$.

(b) Multiply the numerator and the denominator by $\sqrt{\sqrt2-1}$:
\[
\dfrac{1}{\sqrt{1+\sqrt2}}=\dfrac{\sqrt{\sqrt2-1}}{\sqrt{\left(\sqrt2+1\right)\left(\sqrt2-1\right)}}=\dfrac{\sqrt{\sqrt2-1}}{1}=\boxed{\sqrt{\sqrt2-1}}.
\]
(This equals the second product in (a), $u$, not the first one.)
